\documentclass[12pt]{article}

%% Basic document formatting
\usepackage{amsmath, amsthm, amssymb, amsfonts}
\usepackage{mathtools}
\usepackage{mathrsfs}
\usepackage{xspace}
\usepackage{thmtools}
\usepackage{graphicx}
\usepackage{tikz-cd}
\usepackage{color}
\usepackage{setspace}
\usepackage{fancyhdr}
\usepackage{titling}
\usepackage{hyperref}
\usepackage[left=0.4in,right=0.4in,top=1in,bottom=1in]{geometry}
\usepackage{float}
\usepackage{tabularx}
\usepackage{booktabs}
\usepackage[utf8]{inputenc}
\usepackage[english]{babel}
\usepackage{framed}
\usepackage[dvipsnames]{xcolor}
\usepackage{environ}
\usepackage{tcolorbox}
\tcbuselibrary{theorems,skins,breakable}

\usepackage{enumitem}
\setlist[enumerate]{leftmargin=*}
\setlist[enumerate,1]{labelindent=\parindent}
\setlist[enumerate,2]{labelindent=0pt}

% Blackboard Bold
\newcommand{\Z}{\mathbb{Z}}                 % integers
\newcommand{\Zpl}{\mathbb{Z}_{+}}           % positive integers
\newcommand{\Znn}{\mathbb{Z}_{\geq 0}}      % nonnegative integers. 
\newcommand{\Q}{\mathbb{Q}}                 % rationals
\newcommand{\Qpl}{\mathbb{Q}_{+}}           % positive Rationals
\newcommand{\R}{\mathbb{R}}                 % reals
\newcommand{\Rpl}{\mathbb{R}_{+}}           % positive Reals
\newcommand{\C}{\mathbb{C}}                 % complex numbers
\newcommand{\F}{\mathbb{F}}                 % field

% Words
\newcommand{\st}{\text{ such that }}
\newcommand{\wrt}{\text{ with respect to }}
\newcommand{\with}{\text{ with }}
\newcommand{\ie}{\text{, i.e. }}
\newcommand*{\ditto}{\text{---}\texttt{"}\text{---}}

% Operators
\newcommand{\abs}[1]{\left|#1\right|}                   % absolute value
\newcommand{\norm}[1]{\abs{\abs{#1}}}
\newcommand{\floor}[1]{\left\lfloor #1 \right\rfloor}   % floor
\newcommand{\ceil}[1]{\left\lceil #1 \right\rceil}      % ceiling
\newcommand{\powset}[1]{\mathscr{P}(#1)}

% Category Theory 
\renewcommand{\hom}{\text{Hom}}
\newcommand{\homspace}[3]{\hom_{#1}(#2, #3)}
\newcommand{\coproduct}[9]{
  \begin{tikzcd}[ampersand replacement=\&]
      \& \& #1 \arrow[dll, bend right, "#6"] \arrow[dl, "#5"] \\
   #4 \& #3 \arrow[l, "#9"] \\
      \& \& #2 \arrow[ull, bend left, "#8"] \arrow[ul, "#7"]
  \end{tikzcd}
}

\newcommand{\product}[9]{ 
  \begin{tikzcd}[ampersand replacement=\&]
      \& \& #3 \\
      #1 \arrow[r, "#7"] \arrow[urr, bend left, "#5"] \arrow[drr, bend right, "#6"] \& #2 \arrow[ur, "#8"] \arrow[dr, "#9"] \\ 
      \& \& #4
  \end{tikzcd}
}


% Algebra 
\newcommand{\Syl}[2]{\text{Syl}_{#1}{(#2)}}           % set of Sylow-p subgroups 
\newcommand{\<}{\langle}                % \<x\>, subgroup generated by x 
\renewcommand{\>}{\rangle}
\renewcommand{\index}[2]{[#1 : #2]}
\newcommand{\id}{\text{id}}             % identity element
\newcommand{\lhdneq}{%
  \mathrel{\ooalign{$\lneq$\cr\raise.22ex\hbox{$\lhd$}\cr}}}
\newcommand{\order}[1]{\text{o}(#1)}    % order of an element
\newcommand{\spec}{\operatorname{Spec}}
\newcommand{\rad}{\operatorname{rad}}   % radical of an ideal 
\newcommand{\sym}[1]{\mathfrak{S}_{#1}}
\newcommand{\aut}[1]{\text{Aut}(#1)} 
\newcommand{\gal}[2]{\text{Gal}(#1 / #2)} 
\newcommand{\inn}[1]{\text{Inn}(#1)} 
\let\oldcong\cong
\let\oldequiv\equiv
\renewcommand{\cong}{\oldequiv}
\renewcommand{\equiv}{\oldcong}

\newcommand{\triangleleftneq}{\mathrel{\ooalign{%
  $\trianglelefteq$\cr
  \hidewidth\raise0.06ex\hbox{$\not\mkern4mu$}\hidewidth\cr
}}}

\makeatletter % cycle
\newcommand{\cyc}[1]{(\mathbf{\cyc@process#1\relax})}
\def\cyc@process#1#2\relax{%
	#1%
	\ifx\relax#2\relax
	\else
		\,\cyc@process#2\relax
	\fi
}
\makeatother

\newcommand{\GL}[2]{\text{GL}_{#1}(#2)}                             % general linear group
\newcommand{\SL}[2]{\text{SL}_{#1}(#2)}                             % special linear group
\newcommand{\gl}[2]{\mathfrak{gl}_{#1}(#2)}
\renewcommand{\sl}[2]{\mathfrak{sl}_{#1}(#2)}
\newcommand{\so}[2]{\mathfrak{so}_{#1}(#2)}
\newcommand{\su}[1]{\mathfrak{su}(#1)}

% Linear Algebra
\newcommand{\bmat}[1]{\begin{bmatrix}#1\end{bmatrix}}   % bracketed matrix
\newcommand{\rank}{\operatorname{rank}}                 % rank
\newcommand{\nullity}{\operatorname{nullity}}           % nullity
\newcommand{\tr}{\operatorname{tr}}

% Combinatorics
\newcommand{\qnum}[1]{\left[#1\right]_q}                % q-analog
\newcommand{\qfact}[1]{\left[#1\right]!_q}              % q-analog factorial 
\newcommand{\qbinom}[2]{\binom{#1}{#2}_q}               % Gaussian binomial coefficient

% Topology/Analysis
\newcommand{\ball}[2]{\text{B}_{#1}(#2)}  % B_r(x): open r-balls around x
\newcommand{\diam}{\text{diam}}           % diamter of a set in metric space 
\newcommand{\lowint}[2]{\underline{\int_{#1}^{#2}}}
\newcommand{\upint}[2]{\overline{\int}_{#1}^{#2}}

% Misc Notation
\newcommand{\oneton}[1]{\{1, \ldots, #1\}}
\newcommand{\defeq}{\vcentcolon=}   % :=
\newcommand{\eqdef}{=\vcentcolon}   % =: 
\renewcommand{\bf}[1]{\textbf{#1}}
\newcommand{\im}{\operatorname{im}}
\newcommand{\fnrest}[2]{\left.#1\right|_{#2}}
\newcommand{\isom}{\overset{\sim}{\rightarrow}} 


\renewcommand{\thesection}{\arabic{section}.}
\renewcommand{\thesubsection}{(\alph{subsection})}

\newtheorem{claim}{Claim}
\newtheorem*{lemma}{Lemma}

%% Headers & title setup
\newcommand{\course}{Math100C}
\newcommand{\myname}{Jay Ser}
\setlength{\headheight}{14.5pt}
\pagestyle{fancy}
\fancyhf{}
\renewcommand{\headrulewidth}{0.4pt}
\lhead{\course}
\rhead{\myname}
\cfoot{\thepage}
\setlength{\droptitle}{-4em} 
\title{\course\ - WK6} 
\author{\myname}
\date{2026.05.08}

\begin{document}
\maketitle
\thispagestyle{fancy}

\newtheorem*{hint}{Hint}

%------------------------------------------------------------------------------%
\section{Factoring \texorpdfstring{$x^3-2$}{x\^{}3-2} over \texorpdfstring{$F = \Q[x]/\<x^3-2\>$}{F}}
 
Define the field $F = \Q[x]/\<x^3-2\>$. Prove that the polynomial $x^3-2$ does \emph{not} factor completely over $F$.
 
\begin{hint}
If it did, then $F$ would contain a primitive cube root of unity.
\end{hint}

\noindent\dotfill 

Let $f = x^3 - 2$. 
Suppose $f$ fully factors over $F$, i.e., all its roots are in $F$. 
Besides $\sqrt[3]{2}$, $\sqrt[3]{2} \zeta_3$ is a root of $f$. 
Since $F \equiv \Q(\sqrt[3]{2})$ is a field, $\sqrt[3]{2}$ is invertible, so if $\sqrt[3]{2}\zeta_3$ is in $F$, so is $\zeta_3$. 
This gives the extensions 
$$F / \Q(\zeta_3) / \Q,$$
resulting in the degree relation 
\begin{equation} \label{eq:1_tower}
  \index{F}{\Q} = \index{F}{\Q(\zeta_3)} \index{\Q(\zeta_3)}{\Q}
\end{equation}
However, the left hand side of \eqref{eq:1_tower} is 3, whereas the right hand side is a multiple of 2 because the degree $\zeta_3$ is 2 over $\Q$. 
Hence \eqref{eq:1_tower} results in a contradiction; 
$f$ does not fully factor over $F$.

 
% -----------------------------------------------------------------------
\section{Field-Theoretic Characterization of Constructible Points}
 
Give (with proof) a field-theoretic characterization of the complex numbers $x+iy$ for which the point $(x,y)$ is straightedge-and-compass constructible starting from $(0,0)$ and $(1,0)$.

\noindent\dotfill 

It is geometrically clear that $x + yi$ is constructible via straightedge-and-compass if and only if the lengths $x$ and $y$ are both constructible as extensions of $\Q$; 
that is, if and only if $x$ and $y$ are both algebraic over $\Q$ with degree a power of two. 

Suppose $a + bi$ and $c + di$ are constructible. 
Taking some field $K$ of degree a power of two over $\Q$ that contains $a, b, c, d$, (such a field clearly exists)  
$a + c, b + d \in K$ means $a - c$ and $b - d$ are both degree a power of two over $\Q$, hence $(a - c) + (b - d)i$ is constructible. 
This demonstrates constructible elements in $\C$ are closed under addition with inverses. 
Similarly, that $ac, bd, bc, ad \in K$ means the point $(a + bi)(c + di)$ is constructible, so constructible points are closed under multiplication.  
To verify that multiplicative inverses of constructible points are also constructible, note that 
conjugates of constructible points are clearly constructible points, and the lengths of constructible points are also constructible.
Hence $(a + bi)^{-1} = \frac{a - bi}{\abs{a + bi}}$ is constructible. 
This shows that the constructible points in $\C$ form a subfield of $\C$. 

 
% -----------------------------------------------------------------------
\section{A Subfield of the Cyclotomic Field and Constructibility of \texorpdfstring{$\cos(2\pi/17)$}{cos(2\unichar{"03C0}/17)}}
 
Let $p$ be a prime and define $\zeta = e^{2\pi i/p} \in \C$. Let $a$ be a positive integer coprime to $p$, and let $d$ be the order of the image of $a$ in $\F_p^\times$. Prove that
\[
  \sum_{i=0}^{d-1} \zeta^{a^i}
\]
generates an extension of $\Q$ of degree $(p-1)/d$. Then deduce from this that $\cos(2\pi/17)$ is constructible.

\noindent\dotfill 

Have mercy, have mercy upon us. 
 
% -----------------------------------------------------------------------
\section{Primitive Elements in Finite Fields}
 
Let $F$ be a finite field of order $p^a$ for some prime $p$. Prove that the following conditions on $\alpha \in F^\times$ are equivalent:
\begin{enumerate}[label=(\roman*)]
  \item The degree of the minimal polynomial of $\alpha$ over $\F_p$ equals $a$.
  \item The order of $\alpha \in F^\times$ does not divide $p^d - 1$ for any proper divisor $d$ of $a$.
\end{enumerate}
Then deduce that $F \cong \F_p[x]/(f(x))$ for some polynomial $f(x)$ over $\F_p$; that is, the primitive element theorem is true for finite fields (but not for arbitrary fields of positive characteristic).

\noindent\dotfill 

The order of $F$ indicates that it is a degree $a$ extension of $\F_p$. 
Suppose first that the degree of the minimal polynomial of $\alpha$ over $\F_p$ is $d \neq \alpha$. 
Then the extensions $F / \F_{p}(\alpha) / \F_p$ and the corresponding relation of degrees $\index{F}{\F_p} = \index{F}{\F_p(\alpha)}\index{\F_p(\alpha)}{\F_p}$ says $d \mid a$. 
Furthermore, $\F_p(\alpha)$ is a finite field of over $p^d$ and its group of units is cyclic of order $p^d - 1$, so $\alpha$ satisfies 
$$\alpha^{p^d - 1} = 1.$$
The relation continues to hold in $F$, so the order of $\alpha$ in $F^\times$ divides $p^d - 1$, where $d$ is a proper divisor of $a$. 
 
Next, suppose the order of $\alpha$ in $F^\times$ divides $p^d - 1$ for some proper divisor $d \mid a$.
Consider 
$$F' \defeq \{z \in F \mid z^{p^d} = z\}.$$
$F'$ is a subfield of $F$ containing $\F_p$. 
Since the order of $\alpha$ in $F^\times$ divides $p^d - 1$, $\alpha^{p^d} = \alpha$, i.e., $\alpha \in F'$. 
By Fermat's Little Theorem, for all $n \in \F_p$, $n^p = n$, so $n \in F'$. 
Next, the fact that the Frobenius map is an automorphism on $F$ says $F'$ is closed under addition/subtraction. 
Simple calculation shows that $F'$ is closed under multiplication and division as well. 
This shows that $F'$ is a subfield of $F$ containing $\alpha$ and $\F_p$. 
Now, since $F'$ is the collection of the roots of the polynomial $x^{p^d} - x$, $F'$ has at most $p^d$ elements. 
$\alpha$, being an element of $F'$, must therefore have a minimal polynomial over $\F_p$ of degree at most $d$; 
namely, the minimal polynomial for $\alpha$ over $\F_p$ is less than $a$. 
 
% -----------------------------------------------------------------------
\section{Automorphisms of Finite Fields and the Frobenius}
 
Let $F$ be a finite field of order $p^a$ for some prime $p$. Prove that the group of field automorphisms of $F$ is generated by the Frobenius map.
 
\begin{hint}
Write $F \equiv \F_p[x]/(f(x))$; argue that Frobenius acts transitively on the roots of $f$; and check that the stabilizer of any root for the action by all field automorphisms is trivial.
\end{hint}

\noindent \dotfill 

Identify $F \equiv \F_p / \<f\>$ where $f$ is a monic irreducible polynomial of degree $a$ with abstract root $\alpha$. 
Since the Frobenius map is an automorphism on $F$,  
$$f(\alpha^p) = f(\alpha)^p = 0,$$ 
shows that $\alpha, \alpha^p, \ldots, \alpha^{p^{a - 1}}$ are all roots of $f$;
in fact, the following shows that they are \textit{distinct} roots of $f$, thus giving all $a$ roots of $f$: 
suppose $a^{p^k} = a^{p^l}$ with $k < l$ and, for contradiction, $l - k < d$. 
Then $a^{p^k(p^{l - k} - 1)} = 1$ shows 
$$o(\alpha)^{\times} \mid p^k(p^{l - k} - 1),$$
where $o(\alpha)^{\times}$ denotes the order of $\alpha$ is the multiplicative group $F^{\times}$. 
But since $\#F^{\times} = p^a - 1$, $o(\alpha)^{\times} \mid p^{\alpha} - 1$, so namely $p \nmid o(\alpha)^{\times}$. 
Since this means $\gcd(o(\alpha)^{\times}, p^k) = 1$, 
$$o(\alpha) \mid p^{l - k} - 1.$$ 
This contradicts the result from question 4; 
if $\alpha$ has order $a$, then the order of $a$ in $F^{\times}$ does not divide $p^d - 1$ for all proper divisors of $a$. 
Hence $l - k \geq d$.
This proves 
$$\alpha, \alpha^p, \ldots, \alpha^{p^{a - 1}}$$ 
give all the distinct roots of $f$; 
the 1st to $(a - 1)$st powers of the Frobenius automorphism act transitively on the roots of $f$. 

In an abuse of notation, now let $\alpha$ be any root of $f$ and denote by $G \defeq \text{Aut}(F)$, $H \defeq G_{\alpha}$. 
Below, I will show that $H$ is trivial. 
Assuming so, the paragraph above and the orbit-stabilizer theorem says that elements of $G$ correspond to the 1st to $(a - 1)$st powers of the Frobenius map, proving that the Frobenius map generates the group of automorphisms of $F$. 
Suppose $\varphi:F \rightarrow F$ is an automorphism of $F$ that fixes $\alpha$. 
Since $\varphi(1) = 1$ by definition, $\varphi$ fixes $\F_p$ as well. 
Then the commuting diagram 
\begin{center}
  \begin{tikzcd} 
    \F_p[x] \arrow[d, "\pi"] \arrow[dr, "\varphi \pi"] \\ 
    F \arrow[r, "\varphi"] & F
  \end{tikzcd}
\end{center}
where $\pi$ is the canonical projection $\F_p[x] \rightarrow \F_p[x] / \<f\> \equiv F$, shows that $\varphi \pi$ is the identity on the constants and sends $x$ to $\alpha$. 
Thus $\varphi \pi$ is a unique substitution map. 
But namely, such a substitution map is exactly the canonical projection $\F_p[x] \rightarrow F$ since $f$ is the minimal polynomial for $\alpha$. 
Hence 
$$\varphi \pi = \id_{F} \pi.$$
Finally, because $\pi$ is a surjection, i.e., epimorphism, $\varphi = \id_{F}$. 
The only automorphism of $F$ that fixes $\alpha$ is the identity; 
the stabilizer of $\alpha$ is trivial, as was to be shown.  
This finishes the problem!

% -----------------------------------------------------------------------
\section{A Computation in \texorpdfstring{$\F_3[x]/\<x^3-x+1\>$}{F\_3[x]/(x\^{}3-x+1)}}
 
Check that $F = \F_3[x]/\<x^3 - x + 1\>$ is a finite field, then factor the polynomial
\[
  f(x) = x^3 + x^2 - x + 1
\]
into irreducibles over $F$.
 
\begin{hint}
For $\alpha \in F$ the image of $x$, first check that $\alpha^2 + 1$ is a root of $f$.
\end{hint}

\noindent\dotfill

$x^3 - x + 1$ is irreducible over $\F_3$: 
because it is cubic, if it is reducible, then it necessarily has a root. 
However, one can easily compute that none of $0, 1$, or $2$ satisfy $x^3 - x + 1$. 
Since $x^3 - x + 1$ is irreducible, one concludes the field $F$ is a three-dimensional vector space over $\F_3$;
namely, $F$ is a finite field of order 27.  

To verify $\alpha^2 + 1$ is a root of $f$, use the identity $\alpha^3 - \alpha + 1 = 0$ in $F$ to calculate  
\begin{align*}
  \alpha^4 & = \alpha^2 - \alpha \\ 
  \alpha^6 & = \alpha^2 - 2\alpha + 1
\end{align*}
Now, 
\newcommand{\myroot}{\alpha^2 + 1}
\begin{align*}
  (\myroot)^3 + (\myroot)^2 - (\myroot) + 1 
  & = \alpha^6 + 1 + \alpha^4 + 2\alpha^2 + 1 - \alpha^2 - 1 + 1 \\ 
  & = \alpha^6 + \alpha^4 + \alpha^2 +  2 \\ 
  & = 0 
\end{align*}
verifies that $\alpha^2 + 1$ is indeed a root of $f$.
Since $f$ is irreducible over $\F_3$, as one can verify by checking $f$ has no roots, and the Frobenius map is an automorphism on $F$ that fixes $\F_3$, 
\begin{align*}
  (\alpha^2 + 1)^3 = \alpha^6 + 1 & = \alpha^2 + \alpha + 2, \\ 
  ((\alpha^2 + 1)^3)^3 = \alpha^6 + \alpha^3 + 2 & = \alpha^2 + 2\alpha + 2
\end{align*}
are both roots of $f$ as well.  
Namely, $\alpha^2 + 1$, $\alpha^2 + \alpha + 2$, and $\alpha^2 + 2\alpha + 2$ are all distinct roots because $1, \alpha, \alpha^2$ form a basis for $F$ over $\F_3$. 
Hence $f$ factors as 
$$(x - (\alpha^2 + 1))(x - (\alpha^2 + \alpha + 2))(x - (\alpha^2 + 2\alpha + 2))$$
over $F$. 
 
% -----------------------------------------------------------------------
\section{An Algebraic Closure of \texorpdfstring{$\F_p$}{F\_p}}
 
Prove that for any prime $p$, there exists a (non-finite) extension $F$ of $\F_p$ such that:
\begin{enumerate}[label=(\roman*)]
  \item every element of $F$ is algebraic over $\F_p$;
  \item \emph{every} polynomial over $\F_p$ factors completely over $F$.
\end{enumerate}
 
\begin{hint}
List the polynomials over $\F_p$ in order, then make a sequence of finite extensions
\[
  \F_p = K_0 \subset K_1 \subset \cdots
\]
with the property that the first $n$ polynomials in your list all factor completely over $K_n$.
\end{hint}
 
\noindent\textit{Optional.} Show that every polynomial over $F$ also factors completely, which is to say $F$ is \emph{algebraically closed}.
 
\begin{hint}
Prove that every polynomial over $F$ divides some polynomial with coefficients in $\F_p$. For instance, this can be done using symmetric polynomials.
\end{hint}

\noindent\dotfill 

Because $\F_p$ is countable and $\F_p[x]$ is generated as a $\F_p$-module by the countably infinite many elements $1, x, x^2, \ldots$, the number of polynomials over $\F_p$ is countably infinite. 
(Arbitrarily) order the elements of $\F_p[x]$ as $(f_1, f_2, f_3, \ldots)$, and define $K_0 \subset K_1 \subset K_2 \subset \ldots$ as given in the hint. 
Now, consider 
$$K \defeq \bigcup_{i = 0}^{\infty} K_i.$$
$K$ is a field:
if $a, b \in K$, let $i$ and $j$ be such that $a \in K_i$ and $b \in K_j$. 
Since the $K_t$'s form a chain, $a$ and $b$ are both elements of $K_{\max\{i, j\}}$, and hence $a - b, ab^{-1} \in K_{max\{i, j\}} \subset K$.  

By definition, every polynomial over $\F_p$ splits in $K$: 
for example, for any $n$, $f_n$ splits in $K_n \subset K$. 
Also, since every element of $K$ comes from some $K_n$, where $K_n$ is a finite extension of $\F_p$, every element of $K$ is algebraic over $\F_p$. 
Hence $K$ satisfies the desired properties. 

 
% -----------------------------------------------------------------------
\section*{Reading (Not a Problem)}
 
Please read the proof of the Fundamental Theorem of Algebra at the end of Chapter~15 of Artin. A different proof will be given later.

\end{document}
